\DSource{0137}{48}
\hypertarget{AB01-PDF0137-P01}{}
\DLine{AB01-PDF0137-body-L001}{}{}{Lunae et aequatio media quinque planetarum, quae est semidiametrus eorum epicycli, cum}
\DLine{AB01-PDF0137-body-L002}{}{}{eius sinus sumatur et eo modo dividatur.}
\hypertarget{AB01-PDF0137-P02}{}
\DLine{AB01-PDF0137-body-L003}{}{}{\Indent Si id calculo reperire vis, gradus epicycli quos ab apogeo Sol, vel Luna, vel planeta}
\DLine{AB01-PDF0137-body-L004}{}{}{percurrerit, scilicet anomaliam, observa; si gradus minus quam $180^\circ$ sunt, iis utere, sed si}
\DLine{AB01-PDF0137-body-L005}{}{}{$180^\circ$ superant, eos de $360^\circ$ deme et residuum adhibe. Graduum alterutra via repertorum, si}
\DLine{AB01-PDF0137-body-L006}{5}{}{minus quam $90^\circ$ sunt, sinum atque cosinum multiplica in semidiametrum epicycli sideris qui}
\DLine{AB01-PDF0137-body-L007}{}{}{est totius aequationis sinus; ambo producta per semidiametrum divide, et quod ex divisione}
\DLine{AB01-PDF0137-body-L008}{}{p. 73.}{cosinus provenerit adde $60^p$ scilicet semidiametro. Quadratum huius summae adde quadrato}
\DLine{AB01-PDF0137-body-L009}{}{}{eius quod e divisione sinus exierit; huius summae radicem inveni. Deinde quod e divisione}
\DLine{AB01-PDF0137-body-L010}{}{}{sinus provenerat, in semidiametrum multiplica, et productum per inventam radicem divide. ---}
\DLine{AB01-PDF0137-body-L011}{10}{}{Si contra gradus adhibendi superant $90^\circ$, de iis $90^\circ$ deme; sinum et cosinum residui in semi-}
\DLine{AB01-PDF0137-body-L012}{}{}{diametrum epicycli multiplica et per semidiametrum divide; quod e sinu provenerit de $60^p$}
\DLine{AB01-PDF0137-body-L013}{}{}{deme; quadratum residui adde quadrato eius quod e cosinu exierit, et huius summae radi-}
\DLine{AB01-PDF0137-body-L014}{}{}{cem inveni. Postea quod e divisione cosinus provenerat in semidiametrum multiplica, et pro-}
\DLine{AB01-PDF0137-body-L015}{}{}{ductum per radicem inventam divide. --- Quod ex alterutra operatione tibi provenerit, in}
\DLine{AB01-PDF0137-body-L016}{15}{}{arcum converte; arcus erit portio graduum ad anomaliam qua usus es pertinens, scilicet ae-}
\DLine{AB01-PDF0137-body-L017}{}{}{quatio stellae ({\itshape f}).}
\hypertarget{AB01-PDF0137-P03}{}
\DLine{AB01-PDF0137-body-L018}{}{}{\Indent Semidiametrus epicycli Solis est $2^p4'45''$ (1); semidiametrus epicycli Lunae $5^p15'$ (2);}
\DLine{AB01-PDF0137-body-L019}{}{}{Saturni $6^p29'50''$ (3); Iovis $11^p30'0''$ (4); Martis $39^p27'22''$ (5); Veneris $43^p9'0''$ (6); Mer-}
\DLine{AB01-PDF0137-body-L020}{}{}{curii $22^p30'30''$ (7). Hoc ex observationibus patuit et cum calculis congruit; idque, volente}
\DLine{AB01-PDF0137-body-L021}{20}{}{Deo, est sinus aequationis mediae ({\itshape g}).}
\vspace{27.6pt}
\hypertarget{AB01-PDF0137-H01}{}
\DLine{AB01-PDF0137-body-L022}{}{}{\CenterLine{CAPUT XXIX.}}
\vspace{13.8pt}
\DLine{AB01-PDF0137-body-L023}{25}{}{\CenterLine{De inaequalitate nychthemerōn et de aliorum in alia conversione.}}
\vspace{13.8pt}
\hypertarget{AB01-PDF0137-P04}{}
\DLine{AB01-PDF0137-body-L024}{}{}{\Indent Dicit [auctor]: Plurimi homines et vulgus nychthemera existimant aequalia esse, et eo-}
\DLine{AB01-PDF0137-body-L025}{}{}{rum quodque 24 horas amplecti; sed hoc haud verum est, nam nychthemerum medium est}
\DLine{AB01-PDF0137-body-L026}{}{p. 74.}{ortus omnium 360 temporum aequatoris ab horizonte vel a circulo meridiano, eo addito quod}
\DLine{AB01-PDF0137-body-L027}{30}{}{ex aequatoris temporibus ascendit cum 59 minutis quae Sol, motu suo medio, in nychthemero}
\DLine{AB01-PDF0137-body-L028}{}{}{percurrit. Nychthemerum autem inaequale est ortus 360 temporum aequatoris, addito motu}
\DLine{AB01-PDF0137-body-L029}{}{}{inaequali Solis in nychthemero, qui necessarie minor aut maior est quam $59'$.}
\hypertarget{AB01-PDF0137-P05}{}
\DLine{AB01-PDF0137-body-L030}{}{}{\Indent Quoniam omni loco initium a circulo horizontis variat prout in illo variant ascensiones}
\DLine{AB01-PDF0137-body-L031}{}{}{signorum, initium autem ab hora meridiana immobile est neque mutatur, ob aequalitatem}
\DLine{AB01-PDF0137-body-L032}{35}{}{ascensionum signorum in circulo meridiano quolibet loco, in calculo stellarum earumque lo-}
\DLine{AB01-PDF0137-body-L033}{}{}{corum initium diei non ponitur ab ortu vel occasu Solis, sed a tempore meridiei aut mediae}
\DLine{AB01-PDF0137-body-L034}{}{}{noctis (8) ({\itshape a}).}
\hypertarget{AB01-PDF0137-P06}{}
\DLine{AB01-PDF0137-body-L035}{}{}{\Indent Item quoniam omnes motus stellarum in tabulis nonnisi ad aequales dies supputantur,}
\DLine{AB01-PDF0137-body-L036}{}{}{differentia inter nychthemera inaequalia et nychthemera media negligitur. In motu Solis et}
\DLine{AB01-PDF0137-body-L037}{40}{}{reliquarum stellarum differentia tanta non est ut errorem sensibilem gignat; sed error mani-}
\par\vspace{6pt}\hrule height.25pt\vspace{5pt}
\noindent\begin{minipage}[t]{.48\textwidth}\vspace{0pt}\fontsize{9}{11.5}\selectfont\setlength{\GutterOffset}{0pt}
\hypertarget{AB01-PDF0137-N01}{}
\DLine{AB01-PDF0137-notes_left-L001}{}{}{\Indent (1) Cod. $5'$ pro $4'$.}
\hypertarget{AB01-PDF0137-N02}{}
\DLine{AB01-PDF0137-notes_left-L002}{}{}{\Indent (2) Gradus in cod. $6^\circ$.}
\hypertarget{AB01-PDF0137-N03}{}
\DLine{AB01-PDF0137-notes_left-L003}{}{}{\Indent (3) Secundae apud Plat. $7''$.}
\hypertarget{AB01-PDF0137-N04}{}
\DLine{AB01-PDF0137-notes_left-L004}{}{}{\Indent (4) Secundae apud Plat. $5''$.}
\end{minipage}
\hfill\begin{minipage}[t]{.48\textwidth}\vspace{0pt}\fontsize{9}{11.5}\selectfont\setlength{\GutterOffset}{0pt}
\hypertarget{AB01-PDF0137-N05}{}
\DLine{AB01-PDF0137-notes_right-L001}{}{}{\Indent (5) Cod. $19^\circ25'22''$; minuta apud Plat. $55'$.}
\hypertarget{AB01-PDF0137-N06}{}
\DLine{AB01-PDF0137-notes_right-L002}{}{}{\Indent (6) Plat. $44^\circ9'5''$.}
\hypertarget{AB01-PDF0137-N07}{}
\DLine{AB01-PDF0137-notes_right-L003}{}{}{\Indent (7) Gradus in cod. $29^\circ$.}
\hypertarget{AB01-PDF0137-N08}{}
\DLine{AB01-PDF0137-notes_right-L004}{}{}{\Indent (8) {\itshape Almag.} III, 8 (ed. Halma, t. I, p. 208).}
\end{minipage}
